\centerline{\bf \Huge Недетские взвешивания}

\begin{flushright}
        \scriptsize{\it Меня создало общение с людьми.\\
                    Они создают меня, я создаю их.\\
                    (с) Рей Аянами.}
\end{flushright}

\problem{1}
Есть 19 монет, одна из которых фальшивая и весит меньше настоящих, а все
настоящие весят одинаково. Имеются также двое старых двухчашечных весов.
Весы ломаются, если одна из чаш перевесит другую (однако, можно успеть
понять, какая именно чаша перевесила). Можно ли за три взвешивания
определить фальшивую монету (возможно, поломав в итоге весы)?


\problem{2}
Дана 61 монета одинакового внешнего вида. Известно, что две из них —
фальшивые, что все настоящие монеты — одинакового веса, и обе фальшивые
монеты одинакового веса (но неизвестно, тяжелее они или легче). Как
можно узнать, тяжелее фальшивые монеты или легче настоящих, за три
взвешивания на двухчашечных весах без гирь?

\problem{3}
Геологи взяли в экспедицию 80 банок консервов, веса которых известны и
различны (имеется список). Через некоторое время надписи на банках
стерлись, и только завхоз знает, что где. Он может всем это доказать
(т.е. обосновать, что в какой банке находится), не вскрывая консервов и
пользуясь только сохранившимся списком и двухчашечными весами без гирь
со стрелкой, показывающей разницу весов на чашах. Докажите, что ему для
этой цели достаточно 4 взвешиваний; недостаточно 3 взвешиваний.



\problem{4}
Известно, что из 32 одинаковых по виду монет есть две фальшивые, которые
отличаются от остальных по весу (настоящие монеты равны друг другу по
весу, и фальшивые монеты равны друг другу по весу). Как разделить все
монеты на две равные по весу кучки, сделав не более 4 взвешиваний на
двухчашечных весах без гирь?

\problem{5}
У царя Гиерона есть 11 металлических слитков, неразличимых на вид: царь
знает, что их веса (в некотором порядке) равны 1 кг, 2 кг, ..., 11 кг. Еще
у него есть мешок, который порвется, если в него положить больше 11 кг.
Архимед узнал веса всех слитков и хочет доказать Гиерону, что первый
слиток имеет вес 1 кг. За один шаг он может загрузить несколько слитков
в мешок и продемонстрировать Гиерону, что мешок не порвется (рвать мешок
нельзя!). За какое наименьшее чило загрузок Архимед может добиться
требуемого?

\newpage

\centerline{\Huge \bf}

\problem{6}
Среди 16 монет есть 8 тяжелых — весом по 11 г, и 8 легких — весом по 10
г, но неизвестно, какие из монет какого веса. Одна из монет — юбилейная.
Как за три взвешивания на двухчашечных весах без гирь узнать, является
юбилейная монета тяжелой или легкой?

\problem{7}
Есть $n$ монет неизвестной массы и весы. Разрешается положить несколько
монет на одну сторону весов и такое же количество монет на другую чашу
весов. Весы либо указывают, что массы двух кучек равны, либо указывают,
какая кучка тяжелее. Докажите, что потребуется по крайней мере $n - 1$
взвешиваний, чтобы определить, все ли монеты имеют одинаковую массу.